Rozwój inżynierii oprogramowania w tym języków programowania oraz skali projektów informatycznych spowodował, że
automatyczna weryfikacja kodu stała się istotnym elementem współczesnych praktyk wytwarzania i modyfikacji kodu.
Ten rozdział poświęcony jest krytycznej analizie ewolucji tych metod - od klasycznego przeglądu kodu polegającego na
wzajemnej ocenie współpracowników, przez klasyczne algorytmy analizy statycznej, do zaawansowanych systemów opartych na
generatywnej sztucznej inteligencji.

W tym rozdziale szczególną uwagę skupiono na zastosowaniu dużych modeli językowych (LLM) w połączeniu ze
statyczną analizą kodu oraz paradygmatem Retrieval-Augmented Generation (RAG). Rozwiązania składające się z tych elementów
są obecnie jednym z najbardziej obiecujących nurtów rozwoju narzędzi do zapewniania jakości kodu na poziomie repozytorium.
Przy zestawieniu klasycznych metod statycznych oraz uczenia maszynowego, udało się zidentyfikować
obszary deficytowe, w których użycie LLM i RAG jest w stanie je uzupełnić.

Głównym celem tej sekcji, poza spisaniem i usystematyzowaniem aktualnego stanu wiedzy naukowej, jest
zdefiniowanie luki badawczej, którą będzie wypełniać niniejszy projekt i praca dyplomowa.
Szczegółowy przegląd obecnych technologii, wiedzy oraz najlepszych praktyk służy jako fundament pod tezę, że
powszechne zastosowanie architektury RAG w procesach przeglądu kodu jest kluczowe dla zwiększenia ich wydajności i
poprawności, zwłaszcza w realiach zwiększającego się poziomu skomplikowania oraz wymaganej wiedzy biznesowej w projektach
informatycznych.